\section*{Hoofdstuk \ref{ch:g9:inside-the-atom} --- In het atoom}

\begin{solution}{exo:g9:inside-the-atom:1}
In de kern: \omterm{def:g9:inside-the-atom:proton}{protonen} (positief) en \omterm{def:g9:inside-the-atom:proton}{neutronen} (geen lading). Eromheen:
\omterm{def:g9:inside-the-atom:nucleus}{elektronen} (negatief).
\end{solution}

\begin{solution}{exo:g9:inside-the-atom:2}
\ce{^{16}_{8}O}: 8 \omterm{def:g9:inside-the-atom:proton}{protonen}, 8 \omterm{def:g9:inside-the-atom:proton}{neutronen}, 8 \omterm{def:g9:inside-the-atom:nucleus}{elektronen}.
\ce{^{27}_{13}Al}: 13, 14, 13. \ce{^{40}_{20}Ca}: 20, 20, 20.
\end{solution}

\begin{solution}{exo:g9:inside-the-atom:3}
Het heeft evenveel \omterm{def:g9:inside-the-atom:nucleus}{elektronen} (elk met lading $-e$) als \omterm{def:g9:inside-the-atom:proton}{protonen} (elk
met lading $+e$): de ladingen heffen elkaar op.
\end{solution}

\begin{solution}{exo:g9:inside-the-atom:4}
$A = 26 + 30 = 56$: \ce{^{56}_{26}Fe}.
\end{solution}

\begin{solution}{exo:g9:inside-the-atom:5}
Ja: beide hebben $Z = 17$, dus beide zijn chloor. Ze hebben een
verschillend \omterm{def:g9:inside-the-atom:atomic-number}{massagetal} (18 en 20 \omterm{def:g9:inside-the-atom:proton}{neutronen}): het zijn \omterm{def:g9:inside-the-atom:isotope}{isotopen}.
\end{solution}

\begin{solution}{exo:g9:inside-the-atom:6}
Nee: ze hebben een verschillend \omterm{def:g9:inside-the-atom:atomic-number}{atoomnummer} (6 en 7), dus het zijn
verschillende elementen, koolstof en stikstof. \omterm{def:g9:inside-the-atom:isotope}{Isotopen} hebben dezelfde
$Z$.
\end{solution}

\begin{solution}{exo:g9:inside-the-atom:7}
$56 \times 1.67 \times 10^{-27} \approx \qty{9.35e-26}{kg}$.
\end{solution}

\begin{solution}{exo:g9:inside-the-atom:8}
Een \omterm{def:g7:atoms-and-molecules:atom}{atoom} is grotendeels lege ruimte: de kern is piepklein, dus de meeste
deeltjes komen onderweg geen kern tegen en gaan recht door.
\end{solution}

\begin{solution}{exo:g9:inside-the-atom:9}
$10^{-10} / 10^{-15} = 10^{5}$: honderdduizend keer kleiner.
\end{solution}

\begin{solution}{exo:g9:inside-the-atom:10}
De chemie hangt af van de \omterm{def:g9:inside-the-atom:nucleus}{elektronen}, en de \omterm{def:g9:inside-the-atom:isotope}{isotopen} van een element
hebben evenveel \omterm{def:g9:inside-the-atom:nucleus}{elektronen} (dezelfde $Z$).
\end{solution}

\begin{solution}{exo:g9:inside-the-atom:11}
$26 \times 9.11 \times 10^{-31} = \qty{2.37e-29}{kg}$. Deel:
$2.37 \times 10^{-29} / 9.35 \times 10^{-26} \approx 2.5 \times
10^{-4}$, ongeveer \qty{0.025}{\percent} van de massa.
\end{solution}

\begin{solution}{exo:g9:inside-the-atom:12}
6915 \omterm{def:g7:atoms-and-molecules:atom}{atomen} koper-63 en 3085 \omterm{def:g7:atoms-and-molecules:atom}{atomen} koper-65. Gemiddelde:
$(6915 \times 63 + 3085 \times 65)/10\,000 = 63.6$ \omterm{def:g9:inside-the-atom:proton}{nucleonen}.
\end{solution}

\begin{solution}{pb:g9:inside-the-atom:1}
\textbf{1.} \ce{^{35}_{17}Cl} en \ce{^{37}_{17}Cl}.

\textbf{2.} Chloor-35: 17 \omterm{def:g9:inside-the-atom:proton}{protonen}, 18 \omterm{def:g9:inside-the-atom:proton}{neutronen}, 17 \omterm{def:g9:inside-the-atom:nucleus}{elektronen}.
Chloor-37: 17 \omterm{def:g9:inside-the-atom:proton}{protonen}, 20 \omterm{def:g9:inside-the-atom:proton}{neutronen}, 17 \omterm{def:g9:inside-the-atom:nucleus}{elektronen}.

\textbf{3.} Beide hebben 17 \omterm{def:g9:inside-the-atom:proton}{protonen}, dus beide zijn het element chloor;
ze verschillen alleen in hun aantal \omterm{def:g9:inside-the-atom:proton}{neutronen}, dus het zijn \omterm{def:g9:inside-the-atom:isotope}{isotopen}.

\textbf{4.} Ja: de chemie hangt af van de \omterm{def:g9:inside-the-atom:nucleus}{elektronen}, en beide hebben er
17.

\textbf{5.} $35 \times 1.67 \times 10^{-27} = \qty{5.845e-26}{kg}$.

\textbf{6.} $17 \times 9.11 \times 10^{-31} = \qty{1.55e-29}{kg}$:
ongeveer 4000 keer minder dan de kern, verwaarloosbaar.

\textbf{7.} $37 \times 1.67 \times 10^{-27} = \qty{6.179e-26}{kg}$.

\textbf{8.} $\qty{1}{g} = \qty{e-3}{kg}$; $10^{-3} / 5.845 \times
10^{-26} \approx 1.71 \times 10^{22}$ \omterm{def:g7:atoms-and-molecules:atom}{atomen}.

\textbf{9.} \qty{75.8}{\percent} en \qty{24.2}{\percent}.

\textbf{10.} $758 \times 5.845 \times 10^{-26} = \qty{4.430e-23}{kg}$;
$242 \times 6.179 \times 10^{-26} = \qty{1.495e-23}{kg}$.

\textbf{11.} $4.430 \times 10^{-23} + 1.495 \times 10^{-23} =
\qty{5.925e-23}{kg}$.

\textbf{12.} $5.925 \times 10^{-23} / 1000 \approx
\qty{5.93e-26}{kg}$: de gemiddelde massa van een chlooratoom.
\end{solution}
